\documentclass{article}
\usepackage[utf8]{inputenc}
\newcommand{\comment}[1]{ }
\usepackage{tabularx}
\usepackage{hyperref}
\usepackage{mathtools}
\usepackage{float}
\usepackage{subcaption}
\usepackage{textcomp}
\usepackage[a4paper, margin={1in}]{geometry}
\usepackage[english]{babel}
\usepackage{graphicx}
\usepackage{amsmath}
\usepackage{amsfonts}
\usepackage{fontspec}
\setmonofont{FreeMono}
\usepackage{minted}
\setminted{breaklines=true}
\newmintinline[lean]{lean4}{bgcolor=white}
\newminted[leancode]{lean4}{fontsize=\footnotesize}
\usemintedstyle{tango}

\frenchspacing
\usepackage{amsthm}
\usepackage{authblk}
\usepackage{biblatex}
\usepackage{mathpartir}
\usepackage{stmaryrd}
\usepackage{amssymb}
\newtheorem{lemma}{Lemma}
\newtheorem{theorem}{Theorem}
\addbibresource{./refs.bib}

\newcommand{\llb}{\llbracket}
\newcommand{\rrb}{\rrbracket}
\newcommand{\la}{\langle}
\newcommand{\ra}{\rangle}
\newcommand{\arr}{\rightarrow}
\newcommand{\sort}[1]{\text{Sort}\ #1}
\newcommand{\barr}{\arr_{\beta}}
\newcommand{\affarr}{\rightarrowtail}
\newcommand{\affetaarr}{\rightarrowtail_{\eta}}
\newcommand{\etaarr}{\rightarrow_{\eta}}

\newcolumntype{Y}{>{\centering\arraybackslash}X}

\title{DCC \(\beta\)-Reduction Rules}
\author[1]{Lexi Aiello}

\begin{document}

\begin{figure}[h]
  \centering
  \resizebox{\textwidth}{!}{
    \begin{tabularx}{\textwidth}{l >{\centering\arraybackslash}X r}
      \inferrule[Pil]
      {x \barr x'}
      {\Pi(x, y) \barr \Pi(x', y)} &
      \multicolumn{1}{c}{\inferrule[Pir]
      {y \barr y'}
      {\Pi(x, y) \barr \Pi(x, y')}}&
      \inferrule[Eql']
      {x \barr x'}
      {x = y \barr x' = y}\\\\
      \inferrule[Eqr']
      {y \barr y'}
      {x = y \barr x = y'}&
      \multicolumn{1}{c}{\inferrule[Bang]
      {x \barr x'}
      {!(x) \barr !(x')}}&
      \inferrule[Map0]
      {x \barr x'}
      {!\la x \ra \barr !\la x' \ra}\\\\
      \inferrule[Map]
      { }
      {!\la x \ra \&(y) \barr !(x y)}&
      \multicolumn{1}{c}{\inferrule[Unmap]
      { }
      {\$\la x \ra \&(y) \barr \&(x y)}}&
      \inferrule[Unbox]
      { }
      {*(A) \&(x) \barr x}\\\\
      \inferrule[Unmap]
      { }
      {*\la x \ra \&(y) \barr f x}&
      \multicolumn{1}{c}{\inferrule[Pair]
      { }
      {\llb B \rrb \la x, y \ra z \barr \llb B \rrb (M x, N x (M x))}}&
      \inferrule[Eval]
      { }
      {\text{Eval}(A, B, x) \llb B \rrb (y, z) \barr x y z}\\\\
      \inferrule[Quot]
      { }
      {`(B, x) y \barr x}&
      \multicolumn{1}{c}{\inferrule[Id]
      { }
      {\text{Id}(A) x \barr x}}&
      \inferrule[Fstc]
      { }
      {\text{Fstc}(B) \llb B \rrb (x, y) \barr x}\\\\
      \inferrule[Sndc]
      { }
      {\text{Sndc}(B) \llb B \rrb (x, y) \barr y}&
      \multicolumn{1}{c}{\inferrule[Comp']
      { }
      {(x \circ' y) z \barr x (y z)}}&
      \inferrule[Comp]
      { }
      {(x \circ y) z \barr x (y z)}\\\\
      \inferrule[Comp2]
      { }
      {(x \circ\circ y) z \barr x (y z)}&
      \multicolumn{1}{c}{\inferrule[Curr]
      { }
      {\Lambda(A, B, x) y z \barr x \llb B \rrb (y, z)}}&
      \inferrule[Curr']
      { }
      {!\Lambda(A, B, x) \&(y) \&(z) \barr x \&(\llb B \rrb (y, z))}\\\\
      \inferrule[Sigl]
      {x \barr x'}
      {\llb B \rrb (x, y) \barr \llb B \rrb (x', y)}&
      \multicolumn{1}{c}{\inferrule[Sigr]
      {y \barr y'}
      {\llb B \rrb (x, y) \barr \llb B \rrb (x, y')}}&
      \inferrule[Pairl]
      {x \barr x'}
      {\llb B \rrb \la x, y \ra \barr \llb B \rrb \la x', y \ra}\\\\
      \inferrule[Pairr]
      {y \barr y'}
      {\llb B \rrb \la x, y \ra \barr \llb B \rrb \la x, y' \ra}&
      \multicolumn{1}{c}{\inferrule[Sigma]
      { }
      {\Sigma(B, C) x \barr \text{Coup}(C x)}}&
      \inferrule[Coup]
      {x \barr x'}
      {\text{Coup}(x) \barr \text{Coup}(x')}\\\\
      \multicolumn{1}{c}{\inferrule[CongrFun]
      {x \barr x'}
      {x y \barr x' y}}&&
      \multicolumn{1}{c}{\inferrule[CongrArg]
      {y \barr y'}
      {x y \barr x y'}}
    \end{tabularx}
  }
  \caption{Where \(\barr\) represents single-step beta-reduction. Unmap is not strictly necessary, but is included for performance optimization purposes in compiled code.}
\end{figure}


\end{document}
